\documentclass[../../../include/open-logic-section]{subfiles}

\begin{document}

\olfileid{sth}{cardinals}{zfc}
\olsection{$\ZFC$: Tonggak Wigati}

Kanthi nambahaké Pengurutan Becik, kita wis tekan
tonggak teoretis kang pungkasan. Saiki kita nduwèni kabèh
aksioma kang dibutuhaké kanggo~$\ZFC$. Luwih rinci:

\begin{defn}
Téori $\ZFC$ nduwèni aksioma-aksioma iki: Ekstensionalitas,
Gabungan, Pasangan, Himpunan Pangkat, Tanpa Wates, Landhesan,
Pengurutan Becik, lan kabèh instansi skema Pamisahan lan
Panggantian. Kanthi tembung liya, $\ZFC$ nambahaké
Pengurutan Becik marang~$\ZF$.
\end{defn}

$\ZFC$ iku cekakan saka téori himpunan
\emph{Zermelo--Fraenkel} kanthi \emph{Pilihan}. Iki bisa katon
rada anèh, amarga aksioma kang lagi kita tambahaké aran
“Pengurutan Becik”, dudu “Pilihan”. Nanging nalika mengko kita
ngrumusaké {Pilihan}, bakal kabukten yèn Pengurutan Becik
ekivalen (relatif marang~$\ZF$) karo Pilihan
(delengen \olref[choice][woproblem]{thmwochoice}). Dadi
milih salah sijiné minangka aksioma “dhasar” ora nduwèni
bédané. Lan jeneng “$\ZFC$” wis lumrah ing pustaka ilmiah.

\end{document}
